%% LyX 2.4.4 created this file.  For more info, see https://www.lyx.org/.
%% Do not edit unless you really know what you are doing.
\documentclass[english]{article}
\usepackage[T1]{fontenc}
\usepackage[latin9]{inputenc}
\usepackage{enumitem}
\usepackage{amsmath}
\usepackage{amsthm}
\usepackage{amssymb}
\usepackage{geometry}
\geometry{verbose,tmargin=1in,bmargin=1in,lmargin=1in,rmargin=1in}

\makeatletter
%%%%%%%%%%%%%%%%%%%%%%%%%%%%%% Textclass specific LaTeX commands.
\numberwithin{equation}{section}
\numberwithin{figure}{section}
\newlength{\lyxlabelwidth}      % auxiliary length 

%%%%%%%%%%%%%%%%%%%%%%%%%%%%%% User specified LaTeX commands.
\usepackage{fancyhdr}
\usepackage{lastpage}
\pagestyle{fancy}
\usepackage{enumitem}
\usepackage{ifthen}
\everymath=\expandafter{\the\everymath\displaystyle}
%\DeclareMathSizes{10}{10}{10}{10}
\usepackage{multicol}

\usepackage{tikz}
\usepackage{tikz-3dplot}
\usetikzlibrary{calc,arrows}

\usetikzlibrary{patterns}
\usetikzlibrary{plotmarks}
\usepackage{pgfplots}
\pgfplotsset{compat=newest}

\usepackage{calc}
\usepackage[nomessages]{fp}

\usepackage{datetime}
% Added by lyx2lyx
\setlength{\parskip}{\smallskipamount}
\setlength{\parindent}{0pt}

\makeatother

\usepackage{babel}
\begin{document}
\renewcommand{\author}{Dr.\,James Ashe}

\newcommand{\classNum}{336}

\newcommand{\classSec}{A}

\newcommand{\examType}{HW 3.1}

\renewcommand{\date}{\formatdate{16}{10}{2013}}

%%First page's header%%%%%%%%%%%%%%%%%%%%%%
\fancypagestyle{firststyle} %%header settings for first pagr
{
	\fancyhead[L]{MTH \classNum{}(\classSec{}) \author{}\\
	\textbf{\examType{}} ~\date{}}
	\fancyhead[C]{}
	\fancyhead[R]{\textbf{Solutions}}
	\setlength{\headsep}{14pt}
}
%%%%%%%%%%%%%%%%%%%%%%%%%%%%%%%%%%%%%%%%%%

%%Next page's header%%%%%%%%%%%%%%%%%%%%%%%
\setlist{nolistsep}
\lhead{\classNum{}(\classSec{})} 
\chead{\textbf{Solutions}} 
\cfoot{\thepage{}~of~\pageref{LastPage}} 
\setlength{\headsep}{5pt} 
%%%%%%%%%%%%%%%%%%%%%%%%%%%%%%%%%%%%%%%%%%%

%%Various settings; see the preamble for other settings%%%%%
%\setenumerate[0]{leftmargin=12pt,itemindent=0pt,labelwidth=0pt}
\setlist{topsep=.5pt}
\renewcommand{\labelenumi}{\textbf{\arabic{enumi}.}}
\renewcommand{\labelenumii}{\textbf{\alph{enumii}.}}
\thispagestyle{firststyle}
%%%%%%%%%%%%%%%%%%%%%%%%%%%%%%%%%%%%%%%%%%

\newcommand{\xmax}{0}
\newcommand{\xmin}{-\xmax}
\newcommand{\xstep}{0}
\newcommand{\ymax}{0}
\newcommand{\ymin}{-\ymax}
\newcommand{\ystep}{0} 
\newcommand{\dspace}{\vspace{1cm}}
\newcommand{\xa}{0}
\newcommand{\xb}{0}
\newcommand{\ya}{1}
\newcommand{\yb}{1}

\small
\begin{enumerate}[resume]
\item \vspace{0cm}

\begin{enumerate}[resume]
\item No. $\mathbf{0}\notin\left\{ x\in\mathbb{R}^{2}:x_{1}+x_{2}=1\right\} $
since $0+0\neq1$.
\item [\textbf{h.}]Yes. We can show this by proving that the set is $\mbox{Span}\left(\left(2,1,1\right),(1,2,-1)\right)$.
$s\begin{bmatrix}\begin{array}{r}
2\\
1\\
1
\end{array}\end{bmatrix}+t\begin{bmatrix}\begin{array}{r}
1\\
2\\
-1
\end{array}\end{bmatrix}$ gives a set of vectors that is $\mbox{Span}\left(\left(2,1,1\right),(1,2,-1)\right)$.
So we need only show that $\begin{bmatrix}\begin{array}{r}
3\\
0\\
1
\end{array}\end{bmatrix}\in\mbox{Span}\left(\left(2,1,1\right),(1,2,-1)\right)$. So let $s=2$ and $t=-1$ which gives $(3,0,1)=2(2,1,1)-1(1,2,-1)$.
Thus the set is $\mbox{Span}\left(\left(2,1,1\right),(1,2,-1)\right)$.\vfill
\end{enumerate}
\item \vspace{0cm}

\begin{enumerate}
\item No. To show this we need to find the rank of the column matrix associated
with the given span. Taking $\mbox{Span}\left(\left(1,1,1\right),(1,2,2)\right)$,
\[
A=\begin{bmatrix}\begin{array}{rr}
1 & 1\\
1 & 2\\
1 & 2
\end{array}\end{bmatrix}\rightsquigarrow\begin{bmatrix}\begin{array}{rr}
1 & 1\\
1 & 2\\
0 & 0
\end{array}\end{bmatrix}
\]
So $\mbox{rank}(A)=2<3$ and $\mbox{Span}\left(\left(1,1,1\right),(1,2,2)\right)\neq\mathbb{R}^{3}$.
\end{enumerate}
\begin{enumerate}[resume]
\item [\textbf{d.}]Yes. As above we find the rank of the associated column
matrix.
\[
A=\begin{bmatrix}\begin{array}{rrr}
1 & 2 & 0\\
0 & 1 & 1\\
-1 & 1 & 5
\end{array}\end{bmatrix}\rightsquigarrow\begin{bmatrix}\begin{array}{rrr}
1 & 0 & 0\\
0 & 1 & 0\\
0 & 0 & 1
\end{array}\end{bmatrix}
\]
So $\mbox{rank}(A)=3$ and $\mbox{Span}\left((1,0,-1),(2,1,1),(0,1,5\right)=\mathbb{R}^{3}.$\vfill
\end{enumerate}
\item [\textbf{5.}]$V=\left\{ \mathbf{x}\in\mathbb{R}^{n}\,:\,A\mathbf{x}=B\mathbf{x}\right\} $

\begin{proof}
\vspace{0cm}

\begin{enumerate}[resume]
\item [i.](Need to show that $\mathbf{0}\in V$.) $\mathbf{0}\in\mathbb{R}^{3}$
and $A\mathbf{0}=0=B\mathbf{0}$ so then $\mathbf{0}\in V$.
\item [ii.](Need to show that if $\mathbf{v}\in V$ and $c\in\mathbb{R}$
then $c\mathbf{v}\in V$.) Let $c\in\mathbb{R}$ and $\mathbf{v}\in V$.
$A(c\mathbf{v})=cA\mathbf{v}=cB\mathbf{v}=B(c\mathbf{v})$ so $c\mathbf{v}\in V$.
\item [iii.](Need to show that if $\mathbf{v},\mathbf{w}\in V$ then $\mathbf{v}+\mathbf{w}\in V$.)
Let $\mathbf{v},\mathbf{w}\in V$. $A(\mathbf{v}+\mathbf{w})=A\mathbf{v}+A\mathbf{w}=B\mathbf{v}+B\mathbf{w}=B(\mathbf{v}+\mathbf{w})$
so $\mathbf{v}+\mathbf{w}\in V$.
\end{enumerate}
\end{proof}
\end{enumerate}
\vfill
\end{document}
